\documentclass[10pt,a4paper]{amsart}
\usepackage[margin=19mm]{geometry}
\usepackage{amsmath,amssymb,mathtools}
\usepackage[scr=rsfs,cal=euler]{mathalfa}
\usepackage{xfrac}
\usepackage[unicode,hidelinks,bookmarks=false]{hyperref}
\newcommand{\ie}{i.e.\ }
\newcommand{\cf}{cf.\ }
\newcommand{\resp}{respectively\ }
\newcommand{\Subsection}{\subsection}
\providecommand{\nobreakdash}{\nobreak\hbox{-}}
\setlength{\emergencystretch}{2em}
\multlinegap0pt
\usepackage{bm}
\usepackage{bbm}
\usepackage{dsfont}
\usepackage{xspace}
\usepackage{cancel}
\usepackage{mathtools}
\usepackage{amsthm}
\makeatletter
\@ifundefined{theorem}{\newtheorem{theorem}{Theorem}}{}
\@ifundefined{lemma}{\newtheorem{lemma}[theorem]{Lemma}}{}
\makeatother
\title[Uniqueness of synchronized stationary equilibria in the Kura]{Uniqueness of synchronized stationary equilibria in the Kuramoto mean field game}
\author{Sebastian Munoz}
\date{Source: arXiv:2605.13783v1 (CC BY 4.0)}
\begin{document}
\maketitle

\noindent\emph{Source and licence.} Uniqueness of synchronized stationary equilibria in the Kuramoto mean field game. Version arXiv:2605.13783v1 (CC BY 4.0), distributed under the Creative Commons Attribution 4.0 International licence, as recorded in the arXiv licence metadata for this exact version and shown on the abstract page https://arxiv.org/abs/2605.13783v1. Full rights notice: licenses/arxiv-2605.13783-CC-BY-4.0.txt.\par\smallskip

\noindent\textit{Editorial scope.} Counted unit: the Lemma whose source label is \texttt{lem:geometric-mean} in the article (source0.tex lines 831--842, source label \texttt{lem:geometric-mean}), delivered together with the article's own complete original proof of it (source0.tex lines 843--961, one \texttt{proof} environment of 119 lines). No mathematics is written, repaired or re-derived: statement and proof are the article's own text line for line. Required context retained uncounted: lem:u-shape (lines 529--538); eq:chord (lines 629--631); lem:Z-bounds (lines 653--663). Every definition, display and label of those lines is retained unchanged. Omitted from this package and not redistributed: 1--528, 539--628, 632--652, 664--830, 962--1473. Nothing inside a retained excerpt is omitted or altered: every retained body is delivered exactly as the article's lines are written, so no omission receipt of any kind accompanies this package. Citations: the retained excerpts carry no citation command at all, so no bibliography entry of the article is transcribed and no reference is redistributed. Reference adaptation: none was needed and none was made; every reference inside the delivered unit resolves inside it, so no unresolved reference or citation is produced. Printed slips kept literal: no printed word, symbol, hypothesis or number of the counted statement or of its proof was altered. Layout and class: the article's own class is \texttt{amsart} with packages including geometry, amsmath, amssymb, amsthm, hyperref; the wrapper is the lane's pinned \texttt{amsart} editorial wrapper at 10pt on A4, which restates the theorem-like environments the retained text needs and adds the article's own macro definitions that the retained text needs (none), with extra wrapper packages (bm, bbm, dsfont, xspace, cancel, mathtools). The delivered numbering is the unit's own. Assets: no retained span contains a figure environment, a graphics-inclusion command, a PDF-page inclusion, a table or a TikZ picture; no image file is copied into this package, the package declares no asset dependency and no image file is redistributed with it. Licence: version arXiv:2605.13783v1 of this article is distributed under the Creative Commons Attribution 4.0 International licence, as recorded in the arXiv licence metadata for this exact version and shown on the abstract page https://arxiv.org/abs/2605.13783v1. A full rights notice is delivered at licenses/arxiv-2605.13783-CC-BY-4.0.txt. Characters: every retained excerpt is pure ASCII, so the delivered engine, XeLaTeX with its default Latin Modern text font, reports no missing character.
\begin{lemma}[Shape of $u$]\label{lem:u-shape}
On $(0,\pi)$,
\begin{equation}\label{eq:shape}
u>0,\qquad u_{xx}\leq 0,\qquad
\Bigl(\frac{u}{\sin x}\Bigr)_x\geq 0,\qquad
|u_x|\leq\frac{u}{\sin x}.
\end{equation}
Moreover, with $\mathcal B:=1+\lambda+u_x+u\cot x$, $\mathcal B_x\leq 0$ on
$(0,\pi)$.
\end{lemma}

\begin{equation}\label{eq:chord}
|u_x|\leq\frac{u}{\sin x} .
\end{equation}

\begin{lemma}[Bounds for $Z$]\label{lem:Z-bounds}
Let
\[
z:=v_{x\zeta},\qquad Z:=\frac{z}{\sin x}.
\]
Then $z>0$ and $Z>0$ on $(0,\pi)$, and
\[
0\leq Z_x\leq uZ\qquad\text{on }(0,\pi).
\]
The upper bound is equivalent to $(fZ)_x\leq 0$.
\end{lemma}

\begin{lemma}[Geometric-mean monotonicity]\label{lem:geometric-mean}
Let
\[
a:=v_\zeta,\qquad
\Delta_0(x):=a(x)-a(0),\qquad \Delta_\pi(x):=a(\pi)-a(x).
\]
Then
\[
H(x):=e^{-v(x)}\frac{\sqrt{\Delta_0(x)\Delta_\pi(x)}}{\sin x}
\]
is nonincreasing on $(0,\pi)$.
\end{lemma}

\begin{proof}
With $z, Z$ as in Lemma~\ref{lem:Z-bounds} and $z=a_x$,
$z,Z>0$ on $(0,\pi)$ and $(e^{-v}Z)_x\leq 0$.

To identify pointwise bounds on $\Delta_0,\Delta_\pi$ that force $H_x\leq 0$, use
$\sin^2 x=(1-\cos x)(1+\cos x)$ to factor
\[
H^2=\Bigl(e^{-v}\,\frac{\Delta_0}{1-\cos x}\Bigr)\Bigl(e^{-v}\,\frac{\Delta_\pi}{1+\cos x}\Bigr),
\]
so that $2\log H$ is the sum of the two logarithms. Using $(\Delta_0)_x=z$
and $(\Delta_\pi)_x=-z$,
\[
\partial_x\log\Bigl(e^{-v}\frac{\Delta_0}{1-\cos x}\Bigr)=-u+\frac z{\Delta_0}-\frac{\sin x}{1-\cos x},\qquad
\partial_x\log\Bigl(e^{-v}\frac{\Delta_\pi}{1+\cos x}\Bigr)=-u-\frac z{\Delta_\pi}+\frac{\sin x}{1+\cos x}.
\]
The trigonometric remainders are half-angle quantities. With
\begin{equation}\label{eq:rt-def}
r:=\cot\tfrac x2=\frac{\sin x}{1-\cos x}=\frac{1+\cos x}{\sin x},\qquad
t:=\tan\tfrac x2=\frac{\sin x}{1+\cos x}=\frac{1-\cos x}{\sin x},
\end{equation}
which satisfy $rt=1$, the log-derivatives simplify to
\[
\partial_x\log\Bigl(e^{-v}\frac{\Delta_0}{1-\cos x}\Bigr)=-u+\frac{z}{\Delta_0}-r,\qquad
\partial_x\log\Bigl(e^{-v}\frac{\Delta_\pi}{1+\cos x}\Bigr)=-u-\frac{z}{\Delta_\pi}+t.
\]
The first is $\leq 0$ iff $z/\Delta_0\leq u+r$, i.e.\ $\Delta_0\geq z/(u+r)$.
The second is automatically $\leq 0$ where $t\leq u$ (both $-u+t$ and
$-z/\Delta_\pi$ are nonpositive), and where $t>u$ it reduces to
$z/\Delta_\pi\geq t-u$, i.e.\ $\Delta_\pi\leq z/(t-u)$. So it suffices to establish
\begin{equation}\label{eq:a-pm-bounds}
\Delta_0(x)\geq\frac{z(x)}{u(x)+r(x)},\qquad
\Delta_\pi(x)\leq\frac{z(x)}{t(x)-u(x)}\ \text{on $\{t>u\}$.}
\end{equation}

We translate these into integrals via the monotonicity of $e^{-v}Z$.
Since $(e^{-v}Z)_x\leq 0$,
\[
Z(y)\geq e^{v(y)-v(x)}Z(x)\text{ for }y\leq x,\qquad
Z(y)\leq e^{v(y)-v(x)}Z(x)\text{ for }y\geq x,
\]
and, since $a_x=z=Z\sin x$, multiplying by $\sin y$ and integrating gives
\begin{equation}\label{eq:a-ZI}
\Delta_0(x)\geq Z(x)\,I_0(x),\qquad \Delta_\pi(x)\leq Z(x)\,I_\pi(x),
\end{equation}
where
\[
I_0(x):=e^{-v(x)}\int_0^x e^{v(y)}\sin y\,dy,\qquad
I_\pi(x):=e^{-v(x)}\int_x^\pi e^{v(y)}\sin y\,dy
\]
solve, by differentiating under the integral and using $v_x=u$, the
first-order ODEs
\begin{equation}\label{eq:I-odes}
I_0'+uI_0=\sin x,\quad I_0(0)=0;\qquad
I_\pi'+uI_\pi=-\sin x,\quad I_\pi(\pi)=0.
\end{equation}
By~\eqref{eq:a-ZI} and $Z=z/\sin x$, the bounds~\eqref{eq:a-pm-bounds} will follow from
\begin{equation}\label{eq:I-bounds}
I_0(x)\geq\frac{\sin x}{u(x)+r(x)},\qquad
I_\pi(x)\leq\frac{\sin x}{t(x)-u(x)}\ \text{on $\{t>u\}$,}
\end{equation}
which we now prove. Both proofs use the derivative identities
\begin{equation}\label{eq:half-angle-diff}
\sin x\cdot r'=-r,\qquad \sin x\cdot t'=t,
\end{equation}
obtained by differentiating either form in~\eqref{eq:rt-def}.

For the lower bound on $I_0$, set
\[
D_0:=\frac{\sin x}{u+r}.
\]
Differentiating
$D_0(u+r)=\sin x$ and using $r'=-r/\sin x$
from~\eqref{eq:half-angle-diff} and $r\sin x=1+\cos x$
from~\eqref{eq:rt-def} to simplify the resulting algebra yields
\begin{equation}\label{eq:D0-ode}
D_0'+uD_0=\sin x-\frac{u+\sin x\,u_x}{(u+r)^2}.
\end{equation}
The chord bound $|u_x|\leq u/\sin x$ from Lemma~\ref{lem:u-shape} gives
$\sin x\,u_x\geq -u$, so the subtracted term in~\eqref{eq:D0-ode} is
nonnegative. Subtracting~\eqref{eq:D0-ode} from~\eqref{eq:I-odes} and
multiplying by $e^v$ gives
\[
\bigl(e^v(I_0-D_0)\bigr)_x=e^v\bigl[(I_0-D_0)'+u(I_0-D_0)\bigr]\geq 0.
\]
As $x\to 0^+$, $u\to 0$ and $r\sim 2/x$, so $D_0\sim x^2/2$. Likewise
$I_0\sim\int_0^x y\,dy=x^2/2$. Hence $e^v(I_0-D_0)$ vanishes at $0$ and
is nondecreasing, giving $I_0\geq D_0$.

For the upper bound on $I_\pi$, first observe that $\{t>u\}$ is an
interval $(x_*,\pi)$ for some $x_*\in[0,\pi)$. Indeed, from $t'=t/\sin x$
in~\eqref{eq:half-angle-diff} and the chord bound~\eqref{eq:chord},
\[
\Bigl(\frac u t\Bigr)_x=\frac{u_x-u/\sin x}{t}\leq 0,
\]
so $u/t$ is nonincreasing on $(0,\pi)$. Since $u(\pi)=0$ and
$t(\pi)=+\infty$, $u/t\to 0$ at $\pi$, so $\{u/t<1\}=\{t>u\}$ is a
right-neighborhood of $\pi$, and it has the form
$\{t>u\}=(x_*,\pi)$. On this interval, set
\[
D_\pi:=\frac{\sin x}{t-u}.
\]
Differentiating $D_\pi(t-u)=\sin x$ and using $\sin x\cdot t'=t$
together with $t\sin x=1-\cos x$ to simplify gives
\[
D_\pi'+uD_\pi=-\sin x+\frac{\sin x\,u_x-u}{(t-u)^2}.
\]
The chord bound~\eqref{eq:chord} gives $\sin x\,u_x\leq u$, so the
last term on the right is nonpositive. Subtracting this from~\eqref{eq:I-odes} and multiplying
by $e^v$ gives
\[
\bigl(e^v(D_\pi-I_\pi)\bigr)_x=e^v\bigl[(D_\pi-I_\pi)'+u(D_\pi-I_\pi)\bigr]\leq 0.
\]
At $x=\pi^-$ both $I_\pi$ and $D_\pi$ tend to $0$ ($I_\pi$ from its
definition, $D_\pi$ because $\sin x\to 0$ while $t-u\to\infty$).
Integrating backward from $\pi$ yields $D_\pi\geq I_\pi$.

This establishes~\eqref{eq:I-bounds}, hence~\eqref{eq:a-pm-bounds},
hence $2(\log H)_x\leq 0$, so $H$ is nonincreasing.
\end{proof}

\end{document}
